% !TeX root = ../papers/nc-p3.tex
\section{四重の固定母線と三点繊維の二項行列}
\label{sec:fourfold-fixed-ruling}
\begin{theorem}\label{thm:fourfold-fixed-ruling}
被約四辺有限正則窓の整反対八次で、余因子のスキーム次数十八なら $HL=6$ である。
\end{theorem}
\begin{proof}[長さ四の母線の一致]
$5\le HL\le6$ である。五なら前定理から核は有理五次曲線である。
正規化の母線束を $\pi:S\to B=\mathbb P^1$、有限入力射を $f$ とし、
入力四空間を $V$、母線の商二束を $F$ とする。$\deg F=8$。
定理\ref{thm:sectional-cofactor-budget}の跡の合成から、底に長さ四の固定因子 $\Theta$ があり
\[
      f_*L(-\pi^*\Theta)\ \subset\ M_D\ \subset\ f_*L .
\]
実際 $a=5-0-3=2$ であり、左端の像は正規化の線束の跡包含の合成が定める。
一般入力平面に制限すれば $S$ 上の切断は $B$ 自身で、有限直像は平面への制限と可換である。
前定理の例外点の正規化局所環 $S_z$ で、規格化した $M$ に含まれる最大の
$S_z$-イデアルは $\mathfrak m$ に等しい。実際 $S_z/\mathfrak m=k[v]/(p_4)$ の
$\langle1,v\rangle$ に非零イデアルは入らない。次数一以下の代表に $v$ の適当な冪を掛けると
二次となり、この包に属さないからである。ゆえに $\Theta$ は平面切断ごとにこの長さ四の因子と
一致し、入力の一点にスキームとして写る。

二つの一般入力平面についてその二点は相異なるよう選べる。実際全ての一般平面が
固定点を含むことはできない。この二点を結ぶ入力直線を $\Lambda=\mathbb P(W)$ とすると、
$\Theta$ 上の母線写像はスキームとしてこの定数 $W$ である。
二つの平面の交点切断が各々定数点として母線上にあり、相異なる二点が同じ二空間を生成するためである。
従って $W$ の二切断は $F(-\Theta)$ に入る。$\Lambda$ の外では上の跡包含は同型になり、
任意の $[p]$ の $F$ での零スキームの長さは $4-\operatorname{rank}N(p)\le3$。
階数一曲線上の印は高々六個、直線 $\Lambda$ 上は高々三個なので、
少なくとも十一個の孤立した三点繊維が $\Lambda$ の外に必要である。

\end{proof}
\begin{proof}[二切断が一般に独立な場合]
$W\otimes\mathcal O_B\to F(-\Theta)$ の行列式の次数は零なので、
一般に独立なら同型である。$h$ を $\Theta$ の二項四次式として基底を取ると、
$\Lambda$ 外の入力点の切断は、$z\in\mathbb P(V/W)$ の各アフィンチャートで
\[
      (A_z+w_1 h,\ B_z+w_2 h)
\]
である。$A_z,B_z$ は $z$ の二座標に一次な二項四次式、$w_1,w_2$ は平面内の座標である。
この二式が長さ三で消える場合、共通三次式で割った残りは独立な二項一次式である。
比例する場合は零因子が長さ四以上となり得ない。
$S_i$ を二項 $i$ 次式空間とする。各三点繊維は次の四行四列の一次行列の射影核の元を与える。
\[
   S_1^2\longrightarrow S_5/hS_1,\qquad
   (l_1,l_2)\longmapsto l_2 A_z-l_1 B_z\pmod {h S_1}.
\]
逆に $l_1,l_2$ が独立な核元から $w_1 l_2-w_2 l_1=-(l_2A_z-l_1B_z)/h$ で
$w_1,w_2$ が一意かつ正則に決まる。独立な一次式は互いに素ゆえ、得た二四次式は
同じ三次式と $l_1,l_2$ の積である。逆の核元も比例を除き一意である。
射影核が正次元ならこの開部分の点は孤立しない。
行列式が恒等零なら余階数一の部分もパラメータ曲線上の局所切断となり孤立しない。
行列式が非零ならその次数四より、孤立繊維数は四以下である。

\end{proof}
\begin{proof}[二切断が従属な場合]
非零入力ベクトルが全母線に入ることはない（その点で $N$ が核曲線全体を消す）。
従って $W$ は大域切断としては独立である。$F$ における像の飽和線束を $P$、
商を $Q$ とすると、$P(-\Theta)$ に独立な二切断があり $\deg P\ge5$。
$Q$ は $V/W$ の二切断から生成され、次数零なら全母線が一つの入力平面に入るので
$1\le\deg Q=8-\deg P\le3$。$\mathbb P(V/W)$ の一点に対応する $Q$ の切断の
零因子を $D_z$ とすれば、$\Lambda$ 外の入力点の零スキームは $D_z$ に含まれる。
ゆえに次数三でなければ三点繊維はない。

次数三なら、この鉛筆は有限平坦な三重射 $c:B\to\mathbb P(V/W)$ で
$D_z=c^{-1}(z)$ を与える（消える切断の直線を用いる）。
$\deg P=5$、$c_*P$ は階数三、Euler数六ゆえ次数三である。
$V_z\subset V$ を $W$ と方向 $z$ で定まる三空間とし、
これらを束にした $V_\bullet$ は行列式 $\mathcal O(-1)$ を持つ。
$V_z$ の切断は $D_z$ 上で $Q$ 成分が零なので、全ての $z$ について
$V_\bullet\to c_*P$ の三行行列が定まる。これは $B$ 上で
$c^*V_\bullet\to F$ の像が $P$ に入ることと有限射の随伴からも得られる。
その射影核の $\Lambda$ 外の点はちょうど三点繊維である。
行列式の次数は $3-(-1)=4$。前段と同じ核の次元の議論で孤立繊維は高々四個である。
どちらも必要な十一個に満たない。
\end{proof}
\path{computations/verify_fourfold_fixed_ruling.py} は四次行列式、三次共通因子の復元、
および三重射の直像次数を検算する。
